\documentclass[tikz, border=6pt]{standalone}
% Shared preamble of the standalone TikZ figures (compiled by scripts/build_figures.py).
\usepackage{amsmath, amssymb}
\usepackage{xeCJK}
\setCJKmainfont[BoldFont={Hiragino Sans W6}]{Hiragino Sans W4}
\usetikzlibrary{arrows.meta, positioning, decorations.pathreplacing, calc, patterns}
\definecolor{rod}{RGB}{70,110,170}
\definecolor{rodb}{RGB}{200,120,50}
\definecolor{okc}{RGB}{60,140,80}
\definecolor{ngc}{RGB}{190,70,60}
\definecolor{boxc}{RGB}{235,240,248}
\tikzset{
  >={Stealth[length=2.4mm]},
  every node/.style={font=\small},
  box/.style={draw, rounded corners=2pt, fill=boxc, align=flush center, inner sep=4pt},
  rodA/.style={draw, fill=rod!35},
  rodB/.style={draw, fill=rodb!35},
  note/.style={font=\footnotesize, align=flush center},
  expl/.style={font=\small, align=flush left},
  brace/.style={decorate, decoration={brace, amplitude=4pt}},
}

\begin{document}
\begin{tikzpicture}
  \tikzset{sys/.style={draw, rounded corners=2pt, fill=boxc, minimum width=1.2cm, minimum height=1.2cm}}
  \node[sys] (a1) at (0,0) {$x$};
  \node[sys] (b1) at (1.35,0) {$Q-x$};
  \node[note] at (0.675,-1.0) {分配 $x$\\(切り離した状態)};
  \begin{scope}[xshift=5.4cm]
    \node[sys, fill=okc!15] (a2) at (0,0) {};
    \node[sys, fill=okc!15] (b2) at (1.2,0) {};
    \draw[very thick, rodb] (0.6,-0.6) -- (0.6,0.6);
    \node[above] at (0.6,0.6) {$J$};
    \node[note] at (0.6,-1.0) {結合系の平衡 $J$\\合計 $Q$ だけで決まる};
  \end{scope}
  \begin{scope}[xshift=10.8cm]
    \node[sys] (a3) at (0,0) {$u$};
    \node[sys] (b3) at (1.35,0) {$Q-u$};
    \node[note] at (0.675,-1.0) {分配 $u$};
  \end{scope}
  \draw[->, thick] (2.25,0) -- node[above, note] {接触させて待つ} (4.6,0);
  \draw[->, thick] (7.4,0) -- node[above, note] {断熱壁で切る} (9.9,0);
  \node[expl, anchor=west, text width=13cm] at (-0.6,-2.1)
    {どの $x$ から始めても同じ $J$ を通るので $\gamma_A(x) + \gamma_B(Q-x) \prec \gamma_A(u) + \gamma_B(Q-u)$。\\
     エントロピー原理から $S_A(x) + S_B(Q-x) \le S_A(u) + S_B(Q-u)$。};
\end{tikzpicture}
\end{document}
